
\documentclass{article}
\usepackage{colortbl}
\usepackage{makecell}
\usepackage{multirow}
\usepackage{supertabular}

\begin{document}

\newcounter{utterance}

\centering \large Interaction Transcript for game `dond', experiment `coop\_it', episode 8 with qwen3{-}next{-}80b{-}a3b{-}instruct{-}t1.0.
\vspace{24pt}

{ \footnotesize  \setcounter{utterance}{1}
\setlength{\tabcolsep}{0pt}
\begin{supertabular}{c@{$\;$}|p{.15\linewidth}@{}p{.15\linewidth}p{.15\linewidth}p{.15\linewidth}p{.15\linewidth}p{.15\linewidth}}
    \# & $\;$A & \multicolumn{4}{c}{Game Master} & $\;\:$B\\
    \hline

    \theutterance \stepcounter{utterance}  
    & & \multicolumn{4}{p{0.6\linewidth}}{
        \cellcolor[rgb]{0.9,0.9,0.9}{
            \makecell[{{p{\linewidth}}}]{
                \texttt{\tiny{[P1$\langle$GM]}}
                \texttt{State giocando a un gioco di negoziazione cooperativa in cui dovete accordarvi su come dividere una serie di oggetti tra voi e un altro giocatore.} \\
\\ 
\texttt{Regole:} \\
\texttt{(a) A Lei e all'altro giocatore viene dato un insieme di oggetti. Ognuno di voi riceve anche una funzione di valore segreta, che rappresenta il valore di ciascun tipo di oggetto.} \\
\texttt{(b) Si scambiano messaggi con l'altro giocatore per concordare chi si aggiudica gli oggetti. Potete inviare un massimo di 5 messaggi ciascuno, oppure terminare in anticipo facendo una proposta segreta in qualsiasi momento.} \\
\texttt{(c) A ciascuno di voi viene chiesto di inviare una proposta segreta indicando gli oggetti che desiderate, formattata tra parentesi quadre come segue: "{[}Proposta: <numero> <nome oggetto>, <numero> <nome oggetto>, <...>{]}".} \\
\texttt{(d) Se le vostre proposte sono complementari, cioè ci sono abbastanza oggetti per soddisfare entrambe le proposte, a ciascun giocatore viene assegnato un punteggio basato sulla somma dei valori degli oggetti ricevuti. In caso contrario, entrambi ricevono zero punti.} \\
\texttt{(e) L'obiettivo è massimizzare la somma del proprio punteggio e di quello dell'altro giocatore.} \\
\\ 
\texttt{Cominciamo.} \\
\\ 
\texttt{L'insieme degli oggetti disponibili è:} \\
\\ 
\texttt{3 forbici, 1 ditale, 1 forchetta} \\
\\ 
\texttt{La funzione valore segreta è:} \\
\\ 
\texttt{forbici: 3, ditale: 0, forchetta: 1} \\
\\ 
\texttt{IMPORTANTE: i vostri messaggi, a meno che non si tratti di una proposta segreta, vengono trasmessi direttamente all'altro giocatore, quindi non includete alcuna risposta alle regole o testo di annuncio del vostro messaggio. Per fare una proposta segreta, utilizzate il formato indicato. Non utilizzare delle parentesi quadre quando si comunica all'altro giocatore, altrimenti verrà interpretata come una proposta segreta.} \\
            }
        }
    }
    & \\ \\

    \theutterance \stepcounter{utterance}  
    & \multicolumn{4}{p{0.6\linewidth}}{
        \cellcolor[rgb]{0.9,0.9,0.9}{
            \makecell[{{p{\linewidth}}}]{
                \texttt{\tiny{[P1$\rangle$GM]}}
                \texttt{Ho capito i valori degli oggetti: ogni forbice vale 3, il ditale vale 0 e la forchetta vale 1.  } \\
\texttt{Poiché il ditale ha valore 0, non mi interessa.  } \\
\texttt{Ho 3 forbici (valore totale 9) e 1 forchetta (valore 1).  } \\
\texttt{Il valore complessivo è 10.  } \\
\\ 
\texttt{Per massimizzare il punteggio complessivo, dobbiamo dividere gli oggetti in modo che entrambi ottengano qualcosa di valore positivo.  } \\
\texttt{Poiché il ditale è inutile, possiamo trascurarlo.  } \\
\\ 
\texttt{Propongo:  } \\
\texttt{Tu prendi 2 forbici e la forchetta (valore 7), io prendo 1 forbice (valore 3).  } \\
\texttt{Così entrambi guadagniamo, e il totale è 10.  } \\
\\ 
\texttt{Se sei d'accordo, posso fare la proposta segreta ora.  } \\
\texttt{Altrimenti, possiamo discutere.  } \\
\\ 
\texttt{Ti propongo:  } \\
\texttt{Tu prendi 2 forbici e la forchetta, io prendo 1 forbice.  } \\
\\ 
\texttt{Se accetti, fammelo sapere e farò la proposta.} \\
            }
        }
    }
    & & \\ \\

    \theutterance \stepcounter{utterance}  
    & & & \multicolumn{2}{p{0.3\linewidth}}{
        \cellcolor[rgb]{0.9,0.9,0.9}{
            \makecell[{{p{\linewidth}}}]{
                \texttt{\tiny{[GM$|$GM]}}
                \texttt{continue} \\
            }
        }
    }
    & & \\ \\

    \theutterance \stepcounter{utterance}  
    & & & \multicolumn{2}{p{0.3\linewidth}}{
        \cellcolor[rgb]{0.9,0.9,0.9}{
            \makecell[{{p{\linewidth}}}]{
                \texttt{\tiny{[GM$|$GM]}}
                \texttt{Ho capito i valori degli oggetti: ogni forbice vale 3, il ditale vale 0 e la forchetta vale 1.  } \\
\texttt{Poiché il ditale ha valore 0, non mi interessa.  } \\
\texttt{Ho 3 forbici (valore totale 9) e 1 forchetta (valore 1).  } \\
\texttt{Il valore complessivo è 10.  } \\
\\ 
\texttt{Per massimizzare il punteggio complessivo, dobbiamo dividere gli oggetti in modo che entrambi ottengano qualcosa di valore positivo.  } \\
\texttt{Poiché il ditale è inutile, possiamo trascurarlo.  } \\
\\ 
\texttt{Propongo:  } \\
\texttt{Tu prendi 2 forbici e la forchetta (valore 7), io prendo 1 forbice (valore 3).  } \\
\texttt{Così entrambi guadagniamo, e il totale è 10.  } \\
\\ 
\texttt{Se sei d'accordo, posso fare la proposta segreta ora.  } \\
\texttt{Altrimenti, possiamo discutere.  } \\
\\ 
\texttt{Ti propongo:  } \\
\texttt{Tu prendi 2 forbici e la forchetta, io prendo 1 forbice.  } \\
\\ 
\texttt{Se accetti, fammelo sapere e farò la proposta.} \\
            }
        }
    }
    & & \\ \\

    \theutterance \stepcounter{utterance}  
    & & \multicolumn{4}{p{0.6\linewidth}}{
        \cellcolor[rgb]{0.9,0.9,0.9}{
            \makecell[{{p{\linewidth}}}]{
                \texttt{\tiny{[GM$\rangle$P2]}}
                \texttt{State giocando a un gioco di negoziazione cooperativa in cui dovete accordarvi su come dividere una serie di oggetti tra voi e un altro giocatore.} \\
\\ 
\texttt{Regole:} \\
\texttt{(a) A Lei e all'altro giocatore viene dato un insieme di oggetti. Ognuno di voi riceve anche una funzione di valore segreta, che rappresenta il valore di ciascun tipo di oggetto.} \\
\texttt{(b) Si scambiano messaggi con l'altro giocatore per concordare chi si aggiudica gli oggetti. Potete inviare un massimo di 5 messaggi ciascuno, oppure terminare in anticipo facendo una proposta segreta in qualsiasi momento.} \\
\texttt{(c) A ciascuno di voi viene chiesto di inviare una proposta segreta indicando gli oggetti che desiderate, formattata tra parentesi quadre come segue: "{[}Proposta: <numero> <nome oggetto>, <numero> <nome oggetto>, <...>{]}".} \\
\texttt{(d) Se le vostre proposte sono complementari, cioè ci sono abbastanza oggetti per soddisfare entrambe le proposte, a ciascun giocatore viene assegnato un punteggio basato sulla somma dei valori degli oggetti ricevuti. In caso contrario, entrambi ricevono zero punti.} \\
\texttt{(e) L'obiettivo è massimizzare la somma del proprio punteggio e di quello dell'altro giocatore.} \\
\\ 
\texttt{Cominciamo.} \\
\\ 
\texttt{L'insieme degli oggetti disponibili è:} \\
\\ 
\texttt{3 forbici, 1 ditale, 1 forchetta} \\
\\ 
\texttt{La funzione valore segreta è:} \\
\\ 
\texttt{forbici: 1, ditale: 6, forchetta: 1} \\
\\ 
\texttt{IMPORTANTE: i vostri messaggi, a meno che non si tratti di una proposta segreta, vengono trasmessi direttamente all'altro giocatore, quindi non includete alcuna risposta alle regole o testo di annuncio del vostro messaggio. Per fare una proposta segreta, utilizzate il formato indicato. Non utilizzare delle parentesi quadre quando si comunica all'altro giocatore, altrimenti verrà interpretata come una proposta segreta.} \\
\\ 
\\ 
\texttt{Ho capito i valori degli oggetti: ogni forbice vale 3, il ditale vale 0 e la forchetta vale 1.  } \\
\texttt{Poiché il ditale ha valore 0, non mi interessa.  } \\
\texttt{Ho 3 forbici (valore totale 9) e 1 forchetta (valore 1).  } \\
\texttt{Il valore complessivo è 10.  } \\
\\ 
\texttt{Per massimizzare il punteggio complessivo, dobbiamo dividere gli oggetti in modo che entrambi ottengano qualcosa di valore positivo.  } \\
\texttt{Poiché il ditale è inutile, possiamo trascurarlo.  } \\
\\ 
\texttt{Propongo:  } \\
\texttt{Tu prendi 2 forbici e la forchetta (valore 7), io prendo 1 forbice (valore 3).  } \\
\texttt{Così entrambi guadagniamo, e il totale è 10.  } \\
\\ 
\texttt{Se sei d'accordo, posso fare la proposta segreta ora.  } \\
\texttt{Altrimenti, possiamo discutere.  } \\
\\ 
\texttt{Ti propongo:  } \\
\texttt{Tu prendi 2 forbici e la forchetta, io prendo 1 forbice.  } \\
\\ 
\texttt{Se accetti, fammelo sapere e farò la proposta.} \\
            }
        }
    }
    & \\ \\

    \theutterance \stepcounter{utterance}  
    & & & \multicolumn{4}{p{0.6\linewidth}}{
        \cellcolor[rgb]{0.9,0.9,0.9}{
            \makecell[{{p{\linewidth}}}]{
                \texttt{\tiny{[GM$\langle$P2]}}
                \texttt{Ti propongo: tu prendi 1 forbice e la forchetta, io prendo 2 forbici.  } \\
\texttt{Così tu hai valore 4, io ho valore 6, totale 10.  } \\
\texttt{Il ditale lo lasciamo da parte, non ha valore.  } \\
\texttt{Accetti? Se sì, io faccio la proposta segreta ora.} \\
            }
        }
    }
     \\ \\

    \theutterance \stepcounter{utterance}  
    & & & \multicolumn{2}{p{0.3\linewidth}}{
        \cellcolor[rgb]{0.9,0.9,0.9}{
            \makecell[{{p{\linewidth}}}]{
                \texttt{\tiny{[GM$|$GM]}}
                \texttt{continue} \\
            }
        }
    }
    & & \\ \\

    \theutterance \stepcounter{utterance}  
    & & & \multicolumn{2}{p{0.3\linewidth}}{
        \cellcolor[rgb]{0.9,0.9,0.9}{
            \makecell[{{p{\linewidth}}}]{
                \texttt{\tiny{[GM$|$GM]}}
                \texttt{Ti propongo: tu prendi 1 forbice e la forchetta, io prendo 2 forbici.  } \\
\texttt{Così tu hai valore 4, io ho valore 6, totale 10.  } \\
\texttt{Il ditale lo lasciamo da parte, non ha valore.  } \\
\texttt{Accetti? Se sì, io faccio la proposta segreta ora.} \\
            }
        }
    }
    & & \\ \\

    \theutterance \stepcounter{utterance}  
    & & \multicolumn{4}{p{0.6\linewidth}}{
        \cellcolor[rgb]{0.9,0.9,0.9}{
            \makecell[{{p{\linewidth}}}]{
                \texttt{\tiny{[P1$\langle$GM]}}
                \texttt{Ti propongo: tu prendi 1 forbice e la forchetta, io prendo 2 forbici.  } \\
\texttt{Così tu hai valore 4, io ho valore 6, totale 10.  } \\
\texttt{Il ditale lo lasciamo da parte, non ha valore.  } \\
\texttt{Accetti? Se sì, io faccio la proposta segreta ora.} \\
            }
        }
    }
    & \\ \\

    \theutterance \stepcounter{utterance}  
    & \multicolumn{4}{p{0.6\linewidth}}{
        \cellcolor[rgb]{0.9,0.9,0.9}{
            \makecell[{{p{\linewidth}}}]{
                \texttt{\tiny{[P1$\rangle$GM]}}
                \texttt{Accetto perfettamente.  } \\
\texttt{La tua proposta è ottimale: totale 10, nessun oggetto sprecato, equilibrio ragionevole.  } \\
\\ 
\texttt{Fai pure la tua proposta segreta.  } \\
\texttt{Io ti seguo con la mia immediatamente dopo.} \\
            }
        }
    }
    & & \\ \\

    \theutterance \stepcounter{utterance}  
    & & & \multicolumn{2}{p{0.3\linewidth}}{
        \cellcolor[rgb]{0.9,0.9,0.9}{
            \makecell[{{p{\linewidth}}}]{
                \texttt{\tiny{[GM$|$GM]}}
                \texttt{continue} \\
            }
        }
    }
    & & \\ \\

    \theutterance \stepcounter{utterance}  
    & & & \multicolumn{2}{p{0.3\linewidth}}{
        \cellcolor[rgb]{0.9,0.9,0.9}{
            \makecell[{{p{\linewidth}}}]{
                \texttt{\tiny{[GM$|$GM]}}
                \texttt{Accetto perfettamente.  } \\
\texttt{La tua proposta è ottimale: totale 10, nessun oggetto sprecato, equilibrio ragionevole.  } \\
\\ 
\texttt{Fai pure la tua proposta segreta.  } \\
\texttt{Io ti seguo con la mia immediatamente dopo.} \\
            }
        }
    }
    & & \\ \\

    \theutterance \stepcounter{utterance}  
    & & \multicolumn{4}{p{0.6\linewidth}}{
        \cellcolor[rgb]{0.9,0.9,0.9}{
            \makecell[{{p{\linewidth}}}]{
                \texttt{\tiny{[GM$\rangle$P2]}}
                \texttt{Accetto perfettamente.  } \\
\texttt{La tua proposta è ottimale: totale 10, nessun oggetto sprecato, equilibrio ragionevole.  } \\
\\ 
\texttt{Fai pure la tua proposta segreta.  } \\
\texttt{Io ti seguo con la mia immediatamente dopo.} \\
            }
        }
    }
    & \\ \\

    \theutterance \stepcounter{utterance}  
    & & & \multicolumn{4}{p{0.6\linewidth}}{
        \cellcolor[rgb]{0.9,0.9,0.9}{
            \makecell[{{p{\linewidth}}}]{
                \texttt{\tiny{[GM$\langle$P2]}}
                \texttt{{[}Proposta: 1 forbice, 1 forchetta{]}} \\
            }
        }
    }
     \\ \\

    \theutterance \stepcounter{utterance}  
    & & & \multicolumn{2}{p{0.3\linewidth}}{
        \cellcolor[rgb]{0.9,0.9,0.9}{
            \makecell[{{p{\linewidth}}}]{
                \texttt{\tiny{[GM$|$GM]}}
                \texttt{proposal} \\
            }
        }
    }
    & & \\ \\

    \theutterance \stepcounter{utterance}  
    & & & \multicolumn{2}{p{0.3\linewidth}}{
        \cellcolor[rgb]{0.9,0.9,0.9}{
            \makecell[{{p{\linewidth}}}]{
                \texttt{\tiny{[GM$|$GM]}}
                [1, 0, 1]
            }
        }
    }
    & & \\ \\

    \theutterance \stepcounter{utterance}  
    & & \multicolumn{4}{p{0.6\linewidth}}{
        \cellcolor[rgb]{0.9,0.9,0.9}{
            \makecell[{{p{\linewidth}}}]{
                \texttt{\tiny{[P1$\langle$GM]}}
                \texttt{L'altro giocatore ha presentato una proposta segreta. Ora tocca a Lei presentare una sua proposta.} \\
            }
        }
    }
    & \\ \\

    \theutterance \stepcounter{utterance}  
    & \multicolumn{4}{p{0.6\linewidth}}{
        \cellcolor[rgb]{0.9,0.9,0.9}{
            \makecell[{{p{\linewidth}}}]{
                \texttt{\tiny{[P1$\rangle$GM]}}
                \texttt{{[}Proposta: 1 forbice, 1 forchetta{]}} \\
            }
        }
    }
    & & \\ \\

    \theutterance \stepcounter{utterance}  
    & & & \multicolumn{2}{p{0.3\linewidth}}{
        \cellcolor[rgb]{0.9,0.9,0.9}{
            \makecell[{{p{\linewidth}}}]{
                \texttt{\tiny{[GM$|$GM]}}
                \texttt{proposal} \\
            }
        }
    }
    & & \\ \\

    \theutterance \stepcounter{utterance}  
    & & & \multicolumn{2}{p{0.3\linewidth}}{
        \cellcolor[rgb]{0.9,0.9,0.9}{
            \makecell[{{p{\linewidth}}}]{
                \texttt{\tiny{[GM$|$GM]}}
                [1, 0, 1]
            }
        }
    }
    & & \\ \\

    \theutterance \stepcounter{utterance}  
    & & & \multicolumn{2}{p{0.3\linewidth}}{
        \cellcolor[rgb]{0.9,0.9,0.9}{
            \makecell[{{p{\linewidth}}}]{
                \texttt{\tiny{[GM$|$GM]}}
                [[1, 0, 1], [1, 0, 1]]
            }
        }
    }
    & & \\ \\

\end{supertabular}
}

\end{document}
